\documentclass{article}
\usepackage{amsmath}
\usepackage{amssymb}
\begin{document}

This note is about using model-generated rollouts to train the reward-conditioned trajectory model
through the interval consistency identity. Generated rollouts are never fitted as data: no likelihood term
is applied to a generated trajectory or to a model-produced reward label. They only supply prefixes,
continuations, and reward queries on which the identity is enforced. All heads remain trainable.

\section*{Setup}

Let $h_i$ be a trajectory prefix and $x_{i+1:j}$ a continuation, giving $h_j$. A block $x_t$ includes the
modeled actions, states, and termination tokens---not just actions.

Write the three predictions as:
\begin{align*}
p_\theta(x \mid h_i) & \quad \text{NOR trajectory prediction} \\
p_\theta(x \mid h_i, r) & \quad \text{reward-conditioned trajectory prediction} \\
q_\theta(r \mid h_i) & \quad \text{NOR reward prediction}
\end{align*}

A generated continuation is drawn by requesting a reward $c$:
\[
x_{i+1:j} \sim G_\theta(\cdot \mid h_i, c).
\]
Here $G_\theta$ is the generation procedure: reward-conditioned actions, NOR dynamics, and learned
termination. The request $c$ shapes which states get visited. It is not a label: nothing in training
asserts that the generated continuation achieved $c$, or any other reward.

\section*{1. The consistency identity}

For a prefix $h_i$, a continuation $x_{i+1:j}$, and any reward $r$, Bayes consistency requires
\[
p_\theta(x_{i+1:j} \mid h_i, r) \cdot q_\theta(r \mid h_i) = p_\theta(x_{i+1:j} \mid h_i) \cdot q_\theta(r \mid h_j).
\]

Its log residual is:
\[
\Delta_{i,j}(r) = \log p_\theta(x_{i+1:j} \mid h_i, r) - \log p_\theta(x_{i+1:j} \mid h_i) + \log q_\theta(r \mid h_i) - \log q_\theta(r \mid h_j).
\]

The update minimizes:
\[
\mathcal{L}_{\mathrm{cons}} = \mathbb{E}_{i < j} \left[\Delta_{i,j}(r)^2\right].
\]

This gives a relationship rather than a target:
\[
\frac{q_\theta(r \mid h_j)}{q_\theta(r \mid h_i)} = \frac{p_\theta(x_{i+1:j} \mid h_i, r)}{p_\theta(x_{i+1:j} \mid h_i)}.
\]

The continuation should increase belief in a reward precisely when it is more likely under that reward
condition. The identity holds for every $r$, not only a reward the trajectory achieved; $r$ is a query.

With full gradients, the loss can change both trajectory predictions and both reward predictions. No
endpoint is designated the correct teacher. For an isolated block, writing $v=\log p_\theta(x\mid h_i,r)$,
\[
\frac{\partial\mathcal L_{\mathrm{cons}}}{\partial v}=2\Delta_{i,j}(r),
\]
so the update can favor or disfavor the sampled continuation under its query reward, depending on the
other predictions. Nothing in the objective makes an observed token more likely merely because it was
observed.

Because the reward-conditioned prediction is normalized over continuations, the identity summed over all
continuations of $h_i$ gives the marginal relation
\[
q_i(r) = \mathbb{E}_{x \sim p(\cdot \mid h_i)}[q_j(r)].
\]
Given the NOR trajectory model and the terminal anchoring of $q$ on recorded outcomes, the identity therefore
determines the reward predictions at every prefix and the reward-conditioned model at every
prefix--reward pair, including pairs the offline data never contains. Enforcing it where those pairs occur
is the purpose of generated rollouts.

\section*{2. Query rewards from the reward prediction}

Two rewards enter: the request $c$ that shapes generation, and the query $r$ in the residual. Either can
use the model's estimate of which outcomes are plausible at the current prefix. Write
\[
q_i(k) = q_\theta(R=r_k \mid h_i), \qquad k=0,\ldots,K-1,
\]
where the bins are ordered by increasing reward. Read this distribution in NOR mode, before inserting a
requested reward. Choose the request once for a generated continuation; the episode reward condition stays
fixed as the prefix grows.

A general upward bias has the form
\[
\rho_i(k)
= \frac{q_i(k)\,w(k)}{\sum_{\ell=0}^{K-1}q_i(\ell)\,w(\ell)},
\qquad 0<w(0)\leq \cdots \leq w(K-1)<\infty.
\]
For example, exponential tilting uses
\[
w(k)=\exp(\beta r_k), \qquad \beta\geq 0,
\]
with representative reward values $r_k$, rather than treating bin indices as evenly spaced rewards.
A reward assigned zero probability by $q_i$ is never requested. Uniform sampling above a recorded outcome
has no corresponding dependence on the model's feasibility estimate. The plan is to increase $\beta$ over
training, so that ambitious queries become common only once the reward predictions and the conditioned
policy are accurate enough to support them.

This is a soft preference, not a feasibility guarantee. A softmax reward head generally assigns positive
probability to every bin, including impossible ones. An upward bias can amplify erroneous tail mass, and
an extremely strong bias can overwhelm the model's preference for plausible outcomes. Moreover, $q_i$
describes outcomes under the learned NOR trajectory distribution; low probability can mean either an
impossible outcome or a possible outcome that this behavior rarely achieves.

Hard conditioning on rewards above the predicted mean is another choice:
\begin{align*}
\mu_i &= \sum_{k=0}^{K-1}q_i(k)r_k, \\
\rho_i^{\mathrm{tail}}(k)
&= \frac{q_i(k)\,\mathbf{1}\{r_k>\mu_i\}}
         {\sum_{\ell=0}^{K-1}q_i(\ell)\,\mathbf{1}\{r_\ell>\mu_i\}}.
\end{align*}
It requires a fallback when the denominator is zero. It also renormalizes the upper tail to total
probability one, even if that tail consists almost entirely of a tiny amount of erroneous probability.
A mild soft bias avoids that forced concentration.

Any request rule gives the generated-continuation distribution
\[
g_\theta(x\mid h_i)
=\sum_{k=0}^{K-1}\rho_i(k)\,G_\theta(x\mid h_i,r_k).
\]

These are proposed rules. The current implementation requests uniformly above the recorded offline outcome
and queries the residual at the rollout's own recorded bin; the model-adaptive sampler has not been
implemented. The implemented augmentation also fits relabelled rollouts as data, which is outside this plan.

\section*{3. Generated rollouts as a proposal distribution}

Let $\nu(h_i,x,r)$ be any sampling distribution over prefixes, continuations, and reward queries for which
the log factors are defined. The consistency objective is
\[
\mathcal L_{\mathrm{cons}}(\nu)
=\mathbb E_{(h_i,x,r)\sim\nu}
  \left[\Delta_{i,j}(r)^2\right].
\]
A coherent model has $\Delta_{i,j}(r)=0$ pointwise. It therefore has zero consistency loss under every
such proposal distribution. Changing $\nu$ changes which disagreements receive optimization effort,
without changing the identity that defines agreement. With incomplete coverage or limited capacity,
the selected proposal still matters to the learned solution.

This is the reason to run consistency on deliberately unusual prefixes and high-reward queries: they
probe relationships among the model's beliefs without becoming demonstrations that every observed token
should be made more likely. The generated distribution $g_\theta$ of section 2 is one such proposal. Its
role is to reach prefixes that the offline data never produces, such as a far start followed by
goal-directed steps, and to pair them with the reward queries the test asks about. Squared log residuals
measure relative disagreement even for rare events, but rare events must still be sampled or deliberately
queried to receive this signal.

\subsection*{A rare reward marginal does not imply an unreliable conditional policy}

The goal is reliable high reward when a high reward is requested. It does not require high reward to
be common under NOR behavior. For an exactly represented joint distribution with a deterministic
trajectory reward, a supported reward condition satisfies
\[
\Pr_{x\sim p(\cdot\mid h_i,r)}\{R(x)=r\}=1
\]
even if $q_i(r)$ is extremely small. Here $R(x)$ describes the ideal joint distribution; no environment
reward evaluator is added to training. This is a statement about exact conditional sampling, not a
guarantee for the implemented generator that combines conditioned actions with NOR dynamics.

Thus consistency can improve the high-reward conditional by coordinating beliefs already learned
from offline data, without making the NOR reward marginal optimistic. In the actual model, reliable
reward conditioning requires learning the missing relationships; their mathematical existence is not
enough.

\subsection*{Failure modes}

Consistency admits unhelpful solutions. If both token modes ignore the reward and the reward
prediction is constant along a trajectory, every interval residual is zero. It can also make several
incorrect predictions agree. The objective alone does not distinguish a real relationship from a mutually
supported model error; the offline likelihood terms and the terminal anchoring of $q$ are what select among
coherent joints.

The training objective is
\[
\mathcal L
=\mathcal L_{\mathrm{offline}}
+\beta\,\mathbb E_{\nu}[\mathcal L_{\mathrm{cons}}],
\]
with $\mathcal L_{\mathrm{offline}}$ the supervised objective on recorded trajectories and their recorded
outcomes, and $\nu$ ranging over offline rows, generated rollouts, or a mixture. Compare proposals at a
shared checkpoint on the same queries; subsequent repeated-update runs can then reveal how the different
generators evolve. Judge the result by high-reward conditional performance in the real environment, used
only for evaluation, rather than by model agreement or imagined reward alone.

\section*{4. Repeated sampling and the stopping property}

Running consistency on the model's own rollouts raises the question of feedback: an example is sampled,
training changes the model, it is sampled again. For an individual interval query $y=(h_i,x,r)$,
\[
\nabla_\theta\mathcal L_{\mathrm{cons}}(y)
=2\Delta_\theta(y)\nabla_\theta\Delta_\theta(y).
\]
At a model that satisfies the identity on every sampled query,
\[
\Delta_\theta(y)=0
\quad\Longrightarrow\quad
\nabla_\theta\mathcal L_{\mathrm{cons}}(y)=0
\quad\text{for each }y.
\]
Consequently its sampled gradient has zero variance there as well as zero mean. A repeatedly sampled,
already-consistent trajectory receives no additional update merely because it appears again. There is no
``sampled again, therefore imitate again'' mechanism, because nothing is imitated.

This does not establish global stability. Away from agreement, oversampling some queries changes the
updates; shared parameters couple queries; offline fitting continues to move the model; and a wrong or
reward-ignoring model can also be consistent. Finite steps on squared log residuals can overshoot. Every
head, including the reward head, remains trainable throughout.

\subsection*{Remark: relation to temporal-difference learning}

Summed over continuations, the identity is the policy-evaluation Bellman equation for the NOR behavior,
distributional over reward bins (section 1). The pointwise identity is stronger than a sampled Bellman
residual: a squared residual on one sampled successor is nonzero at the true model because $q_i$ is an
expectation, whereas the conditioned prediction absorbs that fluctuation and $\Delta_{i,j}(r)=0$ holds for
every continuation. It also needs no importance correction under an off-policy proposal.

With full gradients the objective is gradient descent on a bounded residual, so it cannot diverge; its
costs are slow propagation and the degenerate agreements above. A stop-gradient on the later reward
endpoint, $b_j$, would turn the update into semi-gradient TD in the log domain: faster and directional from
the terminal anchor, but off-policy bootstrapping with function approximation, with no convergence
guarantee. It is well defined for the local one-step loss, not under the all-interval variance shortcut.
It is not used; it is a possible ablation if propagation proves too slow.

\end{document}
